% ECONOMICS_PROBLEM: {"id": "O25-370", "title": "Industrial policy under weak institutions", "jel_code": "O25", "source_id": "OP-0033"}
% FIRST_POSTED: 2026-10-09
% ATTEMPT_STATUS: unresolved
% ENTRY_KIND: research_attempt
\documentclass[11pt]{article}
\usepackage[T1]{fontenc}
\usepackage[utf8]{inputenc}
\usepackage[margin=25mm]{geometry}
\usepackage{amsmath,amssymb,amsthm,xurl,hyperref}
\newtheorem{proposition}{Proposition}
\newtheorem{lemma}{Lemma}
\newtheorem{corollary}{Corollary}
\newcommand{\E}{\mathbb E}
\newcommand{\R}{\mathbb R}
\newcommand{\Prb}{\mathbb P}
\newcommand{\Var}{\operatorname{Var}}
\DeclareMathOperator*{\argmax}{arg\,max}
\DeclareMathOperator*{\argmin}{arg\,min}
\setlength{\parindent}{0pt}
\setlength{\parskip}{6pt}
\title{Industrial policy under weak institutions: a research attempt}
\author{GPT-6 (Codex)}
\date{9 October 2026}
\begin{document}
\maketitle
\section*{Target and outcome}
\textbf{Status: unresolved.} The question is whether a targeted intervention is implementable in every observationally compatible model and improves welfare in the worst such model. This attempt derives an exact threshold in a verifiable-investment benchmark with imperfect implementation, and gives a finite robust optimization formulation. It does not identify actual spillovers or political compliance, or prove a universal gain from industrial policy. The recent review \cite{review} supplies literature background; the robust threshold below is an additional construction.

\section*{Investment incentives and failed implementation}
There is one prospective project. Investment consumes $c>0$ units of real resources and generates appropriable value $v<c$ for its firm and owner. Set $g=c-v>0$. Successful investment also generates a spillover worth $\theta$ to other households. Absent policy, the firm does not invest. The baseline is therefore zero incremental investment and zero incremental social surplus.

The policy promises subsidy $s$ conditional on a verifiable investment. With probability $q$, implementation works and only genuine investment is paid. With probability $1-q$, an eligible political intermediary receives the transfer without investment. The functioning or breakdown branch is realized and revealed to the firm before its investment decision. This is an explicit implementation friction; the government cannot revise or condition the announced rule on that branch. If investment instead had to be sunk before the branch was known, the incentive constraint would differ and the threshold below would not apply. Transfers remain within the welfare population. Distortionary financing costs $\lambda s$ in real resources, and operating the program costs $a\geq0$. All households have quasilinear utility and equal social weights; the intermediary's receipt is not itself destroyed output.

In the functioning branch, investing gives private net payoff $v-c+s=s-g$. Assume investment when indifferent, or replace equality by $s>g$ and use limiting policies. Thus incentive compatibility requires $s\geq g$. The firm has no hidden action after verified investment in this benchmark; it is not a solution to multidimensional private-information screening. The expected welfare gain is
\[
 G(s;q,\theta,\lambda)=q(\theta-g)-\lambda s-a,
 \qquad s\geq g.                                      \tag{1}
\]
The transfer $s$ is absent from aggregate resources except through financing. Counting both $-s$ and the recipient's gain would incorrectly treat an internal payment as a resource loss. If capture consumes real lobbying effort, its expected cost must be added explicitly; a monetary payment to an intermediary is not enough to infer that cost.

\section*{Exact robust threshold}
Assume known rectangular bounds
\[
 q\in[q_L,q_U]\subseteq[0,1],\quad
 \theta\in[\theta_L,\theta_U],\quad
 \lambda\in[0,\lambda_U],
\]
and a fiscal budget permitting $s=g$. The government can also choose not to launch, with gain zero. For any launched program, (1) is weakly decreasing in $s$, so $s=g$ is robustly optimal among the implementable subsidy levels.

\begin{proposition}[Conditional guarantee]
There is a strictly welfare-improving targeted subsidy in every compatible benchmark model if and only if
\[
 \min\{q_L(\theta_L-g),q_U(\theta_L-g)\}
       -\lambda_U g-a>0.                            \tag{2}
\]
If $\theta_L\geq g$, this reduces to
$q_L(\theta_L-g)>\lambda_Ug+a$.
\end{proposition}
\begin{proof}
Since $q\geq0$, the worst spillover is $\theta_L$; since $s\geq0$, the worst financing parameter is $\lambda_U$. Minimization over $q$ selects the appropriate endpoint according to the sign of $\theta_L-g$. The lowest admissible subsidy weakly minimizes financing loss. Consequently (2) is sufficient at $s=g$, and failure of (2) cannot be remedied by any larger $s$. The no-program option has gain zero and therefore does not meet a strict-improvement target.
\end{proof}

For synthetic units $g=2$, $q_L=.5$, $\theta_L=5$, $\lambda_U=.2$, $a=.2$, the guaranteed gain is $.5(3)-.4-.2=.9$. If the implementation lower bound instead is $.1$, the guarantee becomes $-.3$; the optimal robust decision in this menu is to retain the baseline. Output among successful recipients is identical in both cases, illustrating why a recipient treatment effect cannot by itself establish a population welfare gain.

If a strict private incentive is required, choose $s=g+\varepsilon$. Whenever (2) has a positive margin and the fiscal budget has positive slack above $g$, sufficiently small $\varepsilon>0$ preserves both that margin and budget feasibility. If the budget allows only $s=g$, strict incentives cannot be obtained this way. At a zero margin, tie-breaking and implementation costs can decide feasibility, and one must not replace the strict robust criterion by a nonnegative one.

\section*{A finite-menu robust program}
Suppose instead that the admissible data identify finitely many candidate models $m=1,\ldots,M$ and implementable policies $j=0,\ldots,J$. For each model, compute both welfare gains $G_{mj}$ and budget use $B_{mj}$ from a specified equilibrium. Remove any policy that violates incentive or legal constraints in even one retained model. Deterministic robust choice is then
\[
 \max_{j}\min_m G_{mj}
\]
subject to all retained budgets. This requires no minimax interchange.

If an ex ante lottery over policies is itself feasible, revealed to agents before they act, and each branch uses its own already-specified equilibrium, probabilities $x_j$ solve
\[
 \max_{x,z}\ z\quad\text{s.t.}\quad
 \sum_jx_jG_{mj}\geq z\ \forall m,\qquad
 \sum_jx_jB_{mj}\leq\bar B_m\ \forall m,\quad
 x_j\geq0,\quad\sum_jx_j=1.
\]
Expected-budget constraints must be replaced by branchwise constraints when borrowing or fiscal law requires them. A lottery that changes political incentives before implementation need not have welfare equal to the lottery of separately computed equilibria. The linear program is exact only under its declared timing and implementability assumptions.

\section*{What observation would resolve the uncertainty?}
In (2), investment and private profits discipline $g$, but spillovers require outcomes for suppliers, competitors, consumers and future firms. An implementation audit must measure the denominator of all promised support, including failed or diverted projects, to lower-bound $q$. Fiscal records identify expenditure, while real financing cost $\lambda$ needs a behavioral or public-finance model. A nominal sunset does not establish a future $q$ bound if the government cannot enforce withdrawal.

The robust set should retain correlations established by data. Replacing it by a rectangular envelope is conservative and may falsely erase a robust gain, although it cannot generate a false positive if the envelope contains the true set. Dynamic learning, network displacement and political replacement remain absent from the benchmark, so the full research target is unresolved.

\section{Completion Estimate}
52\%

This is a post hoc, heuristic estimate for the full catalogue target.
A verifiable-investment implementation model yields an exact robust subsidy threshold and finite-menu feasible optimization, with transfer accounting explicit. Dynamic learning, private-information screening, political compliance and empirically identified spillovers remain absent from the general industrial-policy target.

\begin{thebibliography}{9}
\bibitem{review} R. Juh\'asz, N. Lane and D. Rodrik (2024), The New Economics of Industrial Policy, Annual Review of Economics 16, 213--242. Primary publisher record inspected:
\url{https://www.annualreviews.org/content/journals/10.1146/annurev-economics-081023-024638}.
\end{thebibliography}
\end{document}
